\lyseis{Β3.11}{Ορθογώνια, ρόμβοι, τετράγωνα}
\sol{1} α) Σωστό. β) Λάθος (πρέπει να έχει και ορθές γωνίες). γ) Σωστό. δ) Λάθος (είναι ίσες, όχι κάθετες). ε) Λάθος, έχει δύο άξονες. στ) Σωστό.
\sol{2} α) ορθογώνιο, τετράγωνο. β) ρόμβος, τετράγωνο. γ) ορθογώνιο, τετράγωνο. δ) ρόμβος, τετράγωνο. ε) ρόμβος, τετράγωνο. στ) τετράγωνο.
\sol{3} Ελεύθερη κατασκευή. Ορθογώνιο: τέσσερα ισοσκελή τρίγωνα. Ρόμβος: τέσσερα ορθογώνια τρίγωνα. Τετράγωνο: τέσσερα ορθογώνια και ισοσκελή. Στον ρόμβο με γωνία $60^\circ$ η μικρή διαγώνιος ($3$ cm) τον χωρίζει σε δύο ισόπλευρα τρίγωνα.
\sol{4} $\hat{\mathrm{Γ}}=70^\circ$ και $\hat{\mathrm{Β}}=\hat{\mathrm{Δ}}=180^\circ-70^\circ=110^\circ$.
\sol{5} α) $2\cdot(12+7)=38$ cm. β) $36:4=9$ cm. γ) $30:4=7{,}5$ cm.
\sol{6} Οι διαγώνιοι είναι ίσες: $\mathrm{ΒΔ}=14$ cm. Διχοτομούνται, άρα $\mathrm{ΟΑ}=\mathrm{ΟΒ}=\mathrm{ΟΓ}=\mathrm{ΟΔ}=7$ cm.
\sol{7} Η διαγώνιος διχοτομεί την $\hat{\mathrm{Α}}$: $\widehat{\mathrm{ΟΑΒ}}=25^\circ$. Οι διαγώνιοι είναι κάθετες, άρα $\widehat{\mathrm{ΟΒΑ}}=180^\circ-90^\circ-25^\circ=65^\circ$.
\sol{8} Χρειάζεται $4\cdot4{,}5=18$ m περίφραξη. Το τετράγωνο έχει $4$ άξονες συμμετρίας.
\sol{9} Όχι. Σε ορθογώνιο οι διαγώνιοι είναι ίσες, αλλά εδώ $130\neq128$.
\sol{10} Η περίμετρος είναι $4\cdot40=160$ cm. Οι μισές διαγώνιοι είναι $32$ cm και $24$ cm, άρα ένα τρίγωνο έχει πλευρές $32$, $24$, $40$ και περίμετρο $96$ cm.
\sol{11} Σχηματίζουν ορθογώνιο $125$ cm επί $25$ cm. Περίμετρος $2\cdot(125+25)=300$ cm.
\sol{12} Δεν έχει δίκιο. Οι διαγώνιοι του ορθογωνίου είναι ίσες και διχοτομούνται, αλλά δεν είναι κάθετες. Στο ορθογώνιο $5$ cm επί $3$ cm, αν σχεδιάσεις τις διαγωνίους, ο μοιρογνωμόνιος δείχνει γωνία διαφορετική από $90^\circ$ (περίπου $62^\circ$ και $118^\circ$).
\sol{13} Ο ρόμβος είναι συμμετρικός ως προς τη διαγώνιο $\mathrm{ΑΓ}$. Αν διπλώσουμε κατά την $\mathrm{ΑΓ}$, το $\mathrm{Β}$ πέφτει πάνω στο $\mathrm{Δ}$. Ο άξονας συμμετρίας είναι κάθετος στο τμήμα που ενώνει δύο συμμετρικά σημεία, άρα $\mathrm{ΑΓ}\perp\mathrm{ΒΔ}$.
\sol{14} Το τρίγωνο $\mathrm{ΑΒΔ}$ έχει $\hat{\mathrm{Α}}=90^\circ$ και $\mathrm{ΑΒ}=\mathrm{ΑΔ}$, άρα είναι ισοσκελές με γωνίες βάσης ίσες. Άρα $\widehat{\mathrm{ΑΒΔ}}=(180^\circ-90^\circ):2=45^\circ$. (Ισοδύναμα, η διαγώνιος διχοτομεί την ορθή γωνία.)
\sol{15} $\widehat{\mathrm{ΔΒΓ}}=90^\circ-30^\circ=60^\circ$. Επειδή $\mathrm{ΟΑ}=\mathrm{ΟΒ}$, είναι $\widehat{\mathrm{ΟΑΒ}}=30^\circ$ και $\widehat{\mathrm{ΑΟΒ}}=180^\circ-30^\circ-30^\circ=120^\circ$. Η $\widehat{\mathrm{ΒΟΓ}}$ είναι παραπληρωματική της $\widehat{\mathrm{ΑΟΒ}}$: $60^\circ$. Το τρίγωνο $\mathrm{ΟΒΓ}$ έχει δύο γωνίες $60^\circ$ και την τρίτη $60^\circ$, άρα είναι ισόπλευρο.
\sol{16} Η πλευρά είναι $52:4=13$ cm. Άρα $3x-2=13$, $3x=15$ και $x=5$.
\sol{17} Ελεύθερη εργασία. Το τετράγωνο έχει $4$ άξονες και το ορθογώνιο $2$.
\sol{18} Ελεύθερη εργασία. Οι διαγώνιοι είναι ίσες στο ορθογώνιο και στο τετράγωνο, όχι στον ρόμβο (αν δεν είναι τετράγωνος).
